\addchap{\lsPrefaceTitle} \label{ch:preface}

This book is not a traditional grammar guide or a collection of lesson plans, but a foundation for understanding language and a catalyst for informed teaching. It doesn't merely complement existing knowledge -- it lays the groundwork for novices to understand the landscape of language, while offering seasoned educators new vantage points from which to view familiar territory.

Each chapter equips you with tools for a different facet of English, moving from concepts to use. We begin by building a clear and solid understanding of familiar labels and categories and by asking what counts as “standard” and “grammatical”. We then map categories (noun, verb, phrase, clause) onto functions (head, modifier, subject, dependent), and revisit lexical categories—especially adverbs and prepositions—to show how a leaner system accounts for the data (including stubborn items like \textit{on}). From there we introduce syntax trees as a professional lens for teachers rather than a student exercise. Subsequent chapters take up the sound system—suprasegmentals and connected-speech processes—followed by vocabulary (frequency, definition, deliberate practice, idiomaticity) and the writing system (phonics and spelling patterns). We then survey key constructions (questions, negation, relatives) to expose shared structures, shift to fluency and flow across skills and their trade-offs with accuracy, and end with discourse organization (coherence, information packaging) and the dynamics of conversation (turn-taking, pragmatics, cultural norms).

This order is intentional: it lays conceptual ground first (what counts as a category, a standard, a judgment), stabilizes the category–function distinction before representation (trees), fills in the interfaces that shape form and processing (sound, lexicon, orthography), and only then generalizes across patterns (constructions) before widening to performance and interaction (fluency, discourse, conversation). In short, we move from foundations to form, from form to patterns, and from patterns to use.

The chapters are more or less modular. If your priority is classroom practice, begin with fluency and discourse, then consult constructions as needed. If you are preparing learners for pronunciation-intensive work, start with the sound system and writing system, then loop back to categories and functions. For grammar-focused courses, take the default path through categories, functions, lexical classes, and trees before dipping into constructions. Cross-references can be used to find light prerequisites; otherwise, feel free to reorder to fit your cohort, time, and aims.


Throughout the book, you'll find that it often prioritizes areas that are typically underrepresented in TESL training. I'll point out common pitfalls and misunderstandings, not to criticize, but to help you avoid the ruts that many language courses get stuck in.

As you explore this linguistic landscape, remember that language is an ever-changing ecosystem. My goal is not always to provide definitive answers, but to equip you with the fundamental knowledge and critical thinking skills to navigate the terrain of English as you encounter it in your teaching journey.

This book maps the terrain; it's up to you to explore it.

\paragraph*{A note about the grammatical framework}

The framework provided here follows that of \textit{The Cambridge grammar of the English language} (\textit{CGEL} \cite{Huddleston2002}) very closely, but has been simplified for accessibility. For instance, I use a two-level noun and noun phrase analysis instead of \textit{CGEL}'s three-level approach, employ a smaller set of functions, and avoid the issue of function fusion. This book focuses on providing the essential elements and ideas that teachers need to embark on their English-teaching journey, rather than attempting to cover all constructions dealt with in \textit{CGEL}. For a more in-depth treatment, consider \textit{A student's introduction to English grammar} \citep{huddleston2022} or \textit{CGEL} itself.

\paragraph*{A note about Canadian English}

While most of the book applies to any English-language teaching situation, the phonology section presents vowels as spoken in southern Ontario, Canada. Many of these will align with General American English, but I include features such as Canadian raising (e.g., \href{https://en.wiktionary.org/wiki/File:En-ca-raising-house.ogg}{/hʌʊs/},\footnote{Click to go to audio recordings on \textit{English Wiktionary}.} as opposed to the British or General American \href{https://en.wiktionary.org/wiki/File:en-us-house-noun.ogg}{/haʊs/}), or indeed any other pronunciation. This choice reflects my own dialect and provides a specific reference point for phonological discussions.